% ============================================================
% Question Bank: ch02-08 Personal Financial Literacy
% Topic Title: Personal Financial Literacy
% CCSS: Supplementary
% Questions: 30 (18 MC + 7 SA + 5 Graphical)
% ============================================================

% --- Multiple Choice Questions (18) ---

\begin{practiceQuestion}{2-8-q01}{mc}
\begin{questionText}
Which type of card lets you spend money you already have in your bank account?
\end{questionText}
\choiceA{Credit card}
\choiceB{Gift card}
\choiceC{Debit card}
\choiceD{Rewards card}
\correctAnswer{C}
\explanation{A debit card takes money directly from your bank account, so you must already have the funds.}
\end{practiceQuestion}

\begin{practiceQuestion}{2-8-q02}{mc}
\begin{questionText}
You charge $\$300$ on a credit card with an $18\%$ annual interest rate. If you pay it off after $1$ year, how much interest do you owe?
\end{questionText}
\choiceA{$\$18$}
\choiceB{$\$36$}
\choiceC{$\$54$}
\choiceD{$\$318$}
\correctAnswer{C}
\explanation{Interest $= 300 \times 0.18 \times 1 = \$54$.}
\end{practiceQuestion}

\begin{practiceQuestion}{2-8-q03}{mc}
\begin{questionText}
A credit union is most similar to which of the following?
\end{questionText}
\choiceA{A department store}
\choiceB{A bank}
\choiceC{A stock market}
\choiceD{A government office}
\correctAnswer{B}
\explanation{A credit union works like a bank (savings, checking, loans) but is owned by its members and often has lower fees.}
\end{practiceQuestion}

\begin{practiceQuestion}{2-8-q04}{mc}
\begin{questionText}
You buy a $\$500$ laptop using a credit card at $12\%$ annual interest. If you pay it all back after $1$ year, what is the total cost?
\end{questionText}
\choiceA{$\$512$}
\choiceB{$\$560$}
\choiceC{$\$600$}
\choiceD{$\$550$}
\correctAnswer{B}
\explanation{Interest $= 500 \times 0.12 \times 1 = \$60$. Total $= 500 + 60 = \$560$.}
\end{practiceQuestion}

\begin{practiceQuestion}{2-8-q05}{mc}
\begin{questionText}
Bank A charges a $\$7$ monthly fee. Bank B charges a $\$3$ monthly fee. How much do you save per year by choosing Bank B?
\end{questionText}
\choiceA{$\$4$}
\choiceB{$\$36$}
\choiceC{$\$48$}
\choiceD{$\$84$}
\correctAnswer{C}
\explanation{Difference per month: $7 - 3 = \$4$. Per year: $4 \times 12 = \$48$.}
\end{practiceQuestion}

\begin{practiceQuestion}{2-8-q06}{mc}
\begin{questionText}
Which statement about a debit card is TRUE?
\end{questionText}
\choiceA{You borrow money and pay it back later}
\choiceB{You must pay interest on every purchase}
\choiceC{You spend money that is already in your account}
\choiceD{You can spend more than your account balance}
\correctAnswer{C}
\explanation{A debit card draws directly from funds you already have in your checking account.}
\end{practiceQuestion}

\begin{practiceQuestion}{2-8-q07}{mc}
\begin{questionText}
You deposit $\$800$ in a savings account that pays $3\%$ annual interest. How much interest do you earn in $1$ year?
\end{questionText}
\choiceA{$\$3$}
\choiceB{$\$24$}
\choiceC{$\$80$}
\choiceD{$\$240$}
\correctAnswer{B}
\explanation{Interest $= 800 \times 0.03 = \$24$.}
\end{practiceQuestion}

\begin{practiceQuestion}{2-8-q08}{mc}
\begin{questionText}
What does the interest rate on a credit card tell you?
\end{questionText}
\choiceA{The maximum amount you can borrow}
\choiceB{The percent charged for borrowing money}
\choiceC{The monthly fee for the card}
\choiceD{The amount you must pay each month}
\correctAnswer{B}
\explanation{The interest rate is the percent the lender charges you for borrowing money.}
\end{practiceQuestion}

\begin{practiceQuestion}{2-8-q09}{mc}
\begin{questionText}
You buy a $\$120$ jacket. Using a debit card costs $\$120$. Using a credit card at $15\%$ annual interest (paid after $1$ year) costs how much total?
\end{questionText}
\choiceA{$\$120$}
\choiceB{$\$135$}
\choiceC{$\$138$}
\choiceD{$\$150$}
\correctAnswer{C}
\explanation{Interest $= 120 \times 0.15 = \$18$. Total $= 120 + 18 = \$138$.}
\end{practiceQuestion}

\begin{practiceQuestion}{2-8-q10}{mc}
\begin{questionText}
Compared to a bank, a credit union typically offers:
\end{questionText}
\choiceA{Higher fees and higher interest rates on loans}
\choiceB{No savings accounts}
\choiceC{Lower fees and is owned by its members}
\choiceD{Government-backed stock investments}
\correctAnswer{C}
\explanation{Credit unions are member-owned and often charge lower fees than banks.}
\end{practiceQuestion}

\begin{practiceQuestion}{2-8-q11}{mc}
\begin{questionText}
Bank A pays $1.5\%$ interest on savings. Bank B pays $4\%$. You deposit $\$2{,}000$ for $1$ year. How much more does Bank B earn you?
\end{questionText}
\choiceA{$\$25$}
\choiceB{$\$50$}
\choiceC{$\$80$}
\choiceD{$\$5$}
\correctAnswer{B}
\explanation{Bank A: $2{,}000 \times 0.015 = \$30$. Bank B: $2{,}000 \times 0.04 = \$80$. Difference: $80 - 30 = \$50$.}
\end{practiceQuestion}

\begin{practiceQuestion}{2-8-q12}{mc}
\begin{questionText}
You owe $\$250$ on a credit card with $20\%$ annual interest. If you pay nothing for $1$ year, how much will you owe?
\end{questionText}
\choiceA{$\$270$}
\choiceB{$\$300$}
\choiceC{$\$275$}
\choiceD{$\$350$}
\correctAnswer{B}
\explanation{Interest $= 250 \times 0.20 = \$50$. Total owed $= 250 + 50 = \$300$.}
\end{practiceQuestion}

\begin{practiceQuestion}{2-8-q13}{mc}
\begin{questionText}
Which is a disadvantage of using a credit card?
\end{questionText}
\choiceA{You can buy things before you have the cash}
\choiceB{It builds your credit history}
\choiceC{You pay interest if you don't pay the full balance on time}
\choiceD{It is accepted at most stores}
\correctAnswer{C}
\explanation{The main disadvantage is paying extra money (interest) if you don't pay the full balance by the due date.}
\end{practiceQuestion}

\begin{practiceQuestion}{2-8-q14}{mc}
\begin{questionText}
A checking account has a $\$6$ monthly fee and no interest. A savings account has no monthly fee and pays $2\%$ interest per year. You deposit $\$500$ for $1$ year. Which option gives you more money?
\end{questionText}
\choiceA{Checking --- you keep $\$500$}
\choiceB{Savings --- you end with $\$510$}
\choiceC{They are equal at $\$500$}
\choiceD{Checking --- you end with $\$510$}
\correctAnswer{B}
\explanation{Checking: $500 - (6 \times 12) = 500 - 72 = \$428$. Savings: $500 + (500 \times 0.02) = \$510$. Savings is better.}
\end{practiceQuestion}

\begin{practiceQuestion}{2-8-q15}{mc}
\begin{questionText}
You charge $\$400$ on a credit card at $10\%$ annual interest. You pay $\$200$ right away and owe $\$200$ for $1$ year. How much total interest do you pay?
\end{questionText}
\choiceA{$\$40$}
\choiceB{$\$20$}
\choiceC{$\$10$}
\choiceD{$\$60$}
\correctAnswer{B}
\explanation{You owe $\$200$ for $1$ year. Interest $= 200 \times 0.10 = \$20$.}
\end{practiceQuestion}

\begin{practiceQuestion}{2-8-q16}{mc}
\begin{questionText}
Which payment method does NOT involve borrowing money?
\end{questionText}
\choiceA{Credit card}
\choiceB{Loan from a bank}
\choiceC{Debit card}
\choiceD{Line of credit}
\correctAnswer{C}
\explanation{A debit card uses money you already have. The other options all involve borrowing money.}
\end{practiceQuestion}

\begin{practiceQuestion}{2-8-q17}{mc}
\begin{questionText}
You put $\$1{,}500$ in a savings account at $4\%$ annual interest for $1$ year. What is your account balance at the end of the year?
\end{questionText}
\choiceA{$\$1{,}504$}
\choiceB{$\$1{,}540$}
\choiceC{$\$1{,}560$}
\choiceD{$\$1{,}600$}
\correctAnswer{C}
\explanation{Interest $= 1{,}500 \times 0.04 = \$60$. Balance $= 1{,}500 + 60 = \$1{,}560$.}
\end{practiceQuestion}

\begin{practiceQuestion}{2-8-q18}{mc}
\begin{questionText}
Why is it important to pay your credit card bill on time?
\end{questionText}
\choiceA{To earn more interest from the credit card company}
\choiceB{To avoid paying extra interest charges}
\choiceC{To increase the card's interest rate}
\choiceD{To close your bank account}
\correctAnswer{B}
\explanation{Paying on time avoids interest charges. Late payments mean you pay more for the same purchase.}
\end{practiceQuestion}

% --- Short Answer Questions (7) ---

\begin{practiceQuestion}{2-8-q19}{sa}
\begin{questionText}
You charge $\$240$ on a credit card at $15\%$ annual interest. How much interest do you owe after $1$ year?
\end{questionText}
\correctAnswer{$\$36$}
\explanation{$I = 240 \times 0.15 \times 1 = \$36$.}
\end{practiceQuestion}

\begin{practiceQuestion}{2-8-q20}{sa}
\begin{questionText}
Bank A has a $\$4$ monthly fee. Bank B has no monthly fee. How much do you save per year at Bank B?
\end{questionText}
\correctAnswer{$\$48$}
\explanation{$4 \times 12 = \$48$ saved per year.}
\end{practiceQuestion}

\begin{practiceQuestion}{2-8-q21}{sa}
\begin{questionText}
You deposit $\$600$ at $5\%$ annual interest for $1$ year. How much interest do you earn?
\end{questionText}
\correctAnswer{$\$30$}
\explanation{$I = 600 \times 0.05 = \$30$.}
\end{practiceQuestion}

\begin{practiceQuestion}{2-8-q22}{sa}
\begin{questionText}
A $\$350$ purchase is paid with a credit card at $16\%$ annual interest. If paid after $1$ year, what is the total cost?
\end{questionText}
\correctAnswer{$\$406$}
\explanation{Interest $= 350 \times 0.16 = \$56$. Total $= 350 + 56 = \$406$.}
\end{practiceQuestion}

\begin{practiceQuestion}{2-8-q23}{sa}
\begin{questionText}
Name one advantage and one disadvantage of using a credit card instead of a debit card.
\end{questionText}
\correctAnswer{Advantage: buy now, pay later; Disadvantage: you may pay interest}
\explanation{Credit cards let you purchase before having the cash, but you owe interest if you don't pay the full balance on time.}
\end{practiceQuestion}

\begin{practiceQuestion}{2-8-q24}{sa}
\begin{questionText}
Bank A pays $3\%$ interest per year. Bank B pays $1\%$ per year. You deposit $\$1{,}000$ in each for $1$ year. How much more interest does Bank A earn?
\end{questionText}
\correctAnswer{$\$20$}
\explanation{Bank A: $1{,}000 \times 0.03 = \$30$. Bank B: $1{,}000 \times 0.01 = \$10$. Difference: $30 - 10 = \$20$.}
\end{practiceQuestion}

\begin{practiceQuestion}{2-8-q25}{sa}
\begin{questionText}
A credit card charges $24\%$ annual interest. You owe $\$500$ for $6$ months. How much interest do you owe?
\end{questionText}
\correctAnswer{$\$60$}
\explanation{$I = 500 \times 0.24 \times 0.5 = \$60$.}
\end{practiceQuestion}

% --- Graphical Questions (3 GMC + 2 GSA) ---

\begin{practiceQuestion}{2-8-q26}{gmc}
\begin{questionText}
The bar graph below compares the total cost of a $\$200$ purchase paid with a debit card vs.\ a credit card (at $18\%$ annual interest, paid after $1$ year). What is the difference in total cost?

\begin{center}
\begin{tikzpicture}[scale=0.7]
  \draw[thick,->] (0,0) -- (0,5.5) node[above] {Cost (\$)};
  \draw[thick,->] (0,0) -- (6,0);
  % Y axis labels
  \foreach \y/\label in {0/0, 1/50, 2/100, 3/150, 4/200, 5/250} {
    \draw (-0.1,\y) -- (0.1,\y);
    \node[left] at (-0.15,\y) {\footnotesize \$\label};
  }
  % Debit bar
  \fill[funBlue!70] (1,0) rectangle (2.5,4);
  \node[below] at (1.75,0) {\footnotesize Debit};
  \node[above] at (1.75,4) {\footnotesize \$200};
  % Credit bar
  \fill[funRed!70] (3.5,0) rectangle (5,4.72);
  \node[below] at (4.25,0) {\footnotesize Credit};
  \node[above] at (4.25,4.72) {\footnotesize \$236};
\end{tikzpicture}
\end{center}
\end{questionText}
\choiceA{$\$18$}
\choiceB{$\$36$}
\choiceC{$\$200$}
\choiceD{$\$236$}
\correctAnswer{B}
\explanation{Read the graph: $\$236 - \$200 = \$36$. The credit card costs $\$36$ more.}
\end{practiceQuestion}

\begin{practiceQuestion}{2-8-q27}{gmc}
\begin{questionText}
The table below shows the monthly fees at four financial institutions. Which institution has the lowest annual fee?

\begin{center}
\begin{tabular}{|l|c|}
\hline
\textbf{Institution} & \textbf{Monthly Fee} \\
\hline
Big Bank & $\$8$ \\
\hline
Town Credit Union & $\$2$ \\
\hline
City Bank & $\$5$ \\
\hline
Online Bank & $\$0$ \\
\hline
\end{tabular}
\end{center}
\end{questionText}
\choiceA{Big Bank}
\choiceB{Town Credit Union}
\choiceC{City Bank}
\choiceD{Online Bank}
\correctAnswer{D}
\explanation{Online Bank charges $\$0$/month, so the annual fee is $\$0$. That is the lowest.}
\end{practiceQuestion}

\begin{practiceQuestion}{2-8-q28}{gmc}
\begin{questionText}
The bar graph shows interest earned on a $\$1{,}000$ deposit at three different banks after $1$ year. Which bank pays the highest interest rate?

\begin{center}
\begin{tikzpicture}[scale=0.7]
  \draw[thick,->] (0,0) -- (0,5) node[above] {Interest (\$)};
  \draw[thick,->] (0,0) -- (8,0);
  \foreach \y/\label in {0/0, 1/10, 2/20, 3/30, 4/40} {
    \draw (-0.1,\y) -- (0.1,\y);
    \node[left] at (-0.15,\y) {\footnotesize \$\label};
  }
  % Bank X
  \fill[funBlue!70] (0.5,0) rectangle (2,2);
  \node[below] at (1.25,0) {\footnotesize Bank X};
  \node[above] at (1.25,2) {\footnotesize \$20};
  % Bank Y
  \fill[funRed!70] (3,0) rectangle (4.5,3.5);
  \node[below] at (3.75,0) {\footnotesize Bank Y};
  \node[above] at (3.75,3.5) {\footnotesize \$35};
  % Bank Z
  \fill[green!50!black] (5.5,0) rectangle (7,1.5);
  \node[below] at (6.25,0) {\footnotesize Bank Z};
  \node[above] at (6.25,1.5) {\footnotesize \$15};
\end{tikzpicture}
\end{center}
\end{questionText}
\choiceA{Bank X}
\choiceB{Bank Y}
\choiceC{Bank Z}
\choiceD{They all pay the same}
\correctAnswer{B}
\explanation{Bank Y earned $\$35$ on $\$1{,}000$, which is $3.5\%$ — the highest of the three.}
\end{practiceQuestion}

\begin{practiceQuestion}{2-8-q29}{gsa}
\begin{questionText}
The table shows two credit card offers. You plan to charge $\$600$ and pay it back after $1$ year. How much less does Card A cost you in total?

\begin{center}
\begin{tabular}{|l|c|c|}
\hline
 & \textbf{Card A} & \textbf{Card B} \\
\hline
Annual Interest Rate & $10\%$ & $22\%$ \\
\hline
Annual Fee & $\$25$ & $\$0$ \\
\hline
\end{tabular}
\end{center}
\end{questionText}
\correctAnswer{$\$47$}
\explanation{Card A: interest $= 600 \times 0.10 = \$60$, plus $\$25$ fee $= \$85$ extra. Card B: interest $= 600 \times 0.22 = \$132$. Difference: $132 - 85 = \$47$. Card A saves $\$47$.}
\end{practiceQuestion}

\begin{practiceQuestion}{2-8-q30}{gsa}
\begin{questionText}
The bar graph shows the annual fees at two banks and the annual interest earned on a $\$2{,}000$ deposit. Which bank is the better deal, and how much more do you gain (or save) per year by choosing it?

\begin{center}
\begin{tikzpicture}[scale=0.6]
  \draw[thick,->] (0,0) -- (0,6) node[above] {\$};
  \draw[thick,->] (0,0) -- (10,0);
  \foreach \y/\label in {0/0, 1/10, 2/20, 3/30, 4/40, 5/50} {
    \draw (-0.1,\y) -- (0.1,\y);
    \node[left] at (-0.15,\y) {\footnotesize \$\label};
  }
  % Bank P - Fee
  \fill[funRed!60] (0.5,0) rectangle (1.8,3.6);
  \node[below] at (1.15,-0.1) {\footnotesize Fee};
  \node[above] at (1.15,3.6) {\footnotesize \$36};
  % Bank P - Interest
  \fill[funBlue!60] (2.2,0) rectangle (3.5,4);
  \node[below] at (2.85,-0.1) {\footnotesize Interest};
  \node[above] at (2.85,4) {\footnotesize \$40};
  \node[below] at (2,-.8) {\footnotesize Bank P};

  % Bank Q - Fee
  \fill[funRed!60] (5.5,0) rectangle (6.8,1.2);
  \node[below] at (6.15,-0.1) {\footnotesize Fee};
  \node[above] at (6.15,1.2) {\footnotesize \$12};
  % Bank Q - Interest
  \fill[funBlue!60] (7.2,0) rectangle (8.5,3);
  \node[below] at (7.85,-0.1) {\footnotesize Interest};
  \node[above] at (7.85,3) {\footnotesize \$30};
  \node[below] at (7,-.8) {\footnotesize Bank Q};
\end{tikzpicture}
\end{center}
\end{questionText}
\correctAnswer{Bank Q by $\$14$}
\explanation{Bank P net: $40 - 36 = \$4$ gain. Bank Q net: $30 - 12 = \$18$ gain. Bank Q is better by $18 - 4 = \$14$.}
\end{practiceQuestion}
